% Источники: raw/IMG_7646.jpeg -- raw/IMG_7652.jpeg.
% Дополнение задачи 3.41: raw/IMG_7614.jpeg (фотография доски).
% В рукописном заголовке номер исправлен, дата записана как 7/10/26.
% Компиляция: tectonic 10.tex
\ifdefined\lessonpreambleloaded
  \expandafter\endinput
\fi
\let\lessonpreambleloaded\empty
\documentclass[a4paper,11pt]{article}
\usepackage{fontspec}
\setmainfont{Times New Roman}
\usepackage[russian]{babel}
\usepackage{amsmath,amssymb}
\usepackage[margin=22mm]{geometry}
\usepackage{indentfirst}
\usepackage{needspace}
\usepackage{tikz}
\usetikzlibrary{patterns}
\usepackage[hidelinks]{hyperref}
\usepackage{fancyhdr}
\pagestyle{fancy}
\fancyhf{}
\renewcommand{\headrulewidth}{0pt}
\fancyfoot[C]{\footnotesize\url{https://github.com/bmstu-software-engineering/cs-s5-tv-26y}}
\fancyfoot[R]{\thepage}
\setlength{\emergencystretch}{2em}
\newcommand{\sourcepage}[1]{\par\Needspace{6\baselineskip}\medskip\noindent{\footnotesize\itshape Источник: \detokenize{#1}.}\par\nopagebreak\smallskip}
\newcommand{\editorial}[1]{\par\smallskip\noindent{\small\itshape Примечание к переносу. #1}\par\smallskip}
\newcommand{\lessontitle}[3]{\begin{center}{\large\bfseries Теория вероятностей. #1}\\[4pt]{\large #2}\\[4pt]#3\end{center}}
\newcommand{\parttitle}[1]{\par\Needspace{4\baselineskip}\medskip\noindent\textbf{#1}\par\nopagebreak\smallskip}
\tikzset{element/.style={draw,rectangle,minimum width=7mm,minimum height=5mm,fill=white},every picture/.style={line width=.5pt}}

\begin{document}
\lessontitle{Семинар №\,10}{Формула Байеса. Схема Бернулли. Формула Пуассона}{7/10/26 (дата в исходнике)}
\editorial{Номер в рукописном заголовке исправлен и читается неуверенно.
Дата сохранена буквально, без согласования с датами предыдущих семинаров.}

\sourcepage{IMG_7646.jpeg}
\section*{Полная вероятность и формула Байеса}
Гипотезы $H_1,\ldots,H_n$ образуют полную группу:
\[
 H_i\cap H_j=\varnothing\quad(i\ne j),\qquad
 H_1\cup\cdots\cup H_n=\Omega,\qquad \sum_{i=1}^nP(H_i)=1.
\]
\[
 P(A)=\sum_{i=1}^nP(H_i)P(A\mid H_i).
\]
Формула Байеса при $P(A)>0$:
\[
 \boxed{P(H_i\mid A)=\frac{P(H_i)P(A\mid H_i)}{P(A)}.}
\]

\section*{Задача 3.32. Два стрелка}
Два стрелка независимо делают по одному выстрелу в мишень.
Вероятности попадания первого и второго равны $0{,}8$ и $0{,}4$.
В мишени обнаружена одна пробоина. Найти вероятность того,
что попал первый стрелок.

$A$ --- попал ровно один стрелок; $B$ --- попал первый,
$C$ --- попал второй. Тогда
\[
 P(B)=0{,}8,\quad P(\overline B)=0{,}2,\quad
 P(C)=0{,}4,\quad P(\overline C)=0{,}6.
\]
Выберем гипотезы:
\begin{align*}
 H_1&=B\cap\overline C,& P(H_1)&=0{,}8\cdot0{,}6=0{,}48,\\
 H_2&=\overline B\cap C,&P(H_2)&=0{,}2\cdot0{,}4=0{,}08,\\
 H_3&=B\cap C,&P(H_3)&=0{,}32,\\
 H_4&=\overline B\cap\overline C,&P(H_4)&=0{,}12.
\end{align*}
Проверка: $0{,}48+0{,}08+0{,}32+0{,}12=1$.

\sourcepage{IMG_7647.jpeg}
\[
 P(A\mid H_1)=P(A\mid H_2)=1,\qquad
 P(A\mid H_3)=P(A\mid H_4)=0.
\]
\[
 P(A)=0{,}48\cdot1+0{,}08\cdot1+0{,}32\cdot0+0{,}12\cdot0=0{,}56.
\]
По формуле Байеса
\[
 P(H_1\mid A)=\frac{P(H_1)P(A\mid H_1)}{P(A)}=\frac67.
\]

\section*{Ошибки отбраковки}
В партии деталей $10\%$ брака. При проверке брак обнаруживается с вероятностью
$0{,}95$, исправная деталь ошибочно отбраковывается с вероятностью $0{,}02$.
Найти вероятность того, что отбракованная деталь была исправной.

$A$ --- деталь отбракована; $H_1$ --- деталь исправна,
$H_2$ --- бракованна:
\[
 P(H_1)=0{,}9,\qquad P(H_2)=0{,}1.
\]
\editorial{В словесном условии вероятность ошибочной отбраковки равна $0{,}02$,
но в решении используется $0{,}03$. Противоречие сохранено явно:
ниже перенесены вычисления именно с $0{,}03$, как на фотографии.}
\[
 P(A\mid H_1)=0{,}03,\qquad P(A\mid H_2)=0{,}95,
\]
\[
 P(A)=P(H_1)P(A\mid H_1)+P(H_2)P(A\mid H_2)
 =0{,}9\cdot0{,}03+0{,}1\cdot0{,}95=0{,}122,
\]
\[
 P(H_1\mid A)=\frac{P(H_1)P(A\mid H_1)}{P(A)}
 =\frac{0{,}9\cdot0{,}03}{0{,}122}.
\]
\editorial{Продолжение последнего числового результата находится за правым краем снимка.}

\sourcepage{IMG_7648.jpeg}
\section*{Схема Бернулли}
Проводятся $n$ независимых испытаний. В каждом из них возможны успех $A$
с вероятностью $p=P(A)$ и неудача $\overline A$ с вероятностью $q=1-p$.
Вероятность ровно $k$ успехов:
\[
 \boxed{P_n(k)=C_n^k p^kq^{n-k}.}
\]
Для числа успехов $K$:
\[
 P(k_1\le K\le k_2)=\sum_{k=k_1}^{k_2}C_n^kp^kq^{n-k}.
\]

\section*{Десять бросков монеты}
Монету бросают $10$ раз. Найти вероятность:
\begin{enumerate}
 \item ровно пяти орлов;
 \item не более пяти орлов;
 \item хотя бы одного орла.
\end{enumerate}
Испытания независимы; $p=q=0{,}5$.
\parttitle{1. Ровно пять орлов}
\[
 P_{10}(5)=C_{10}^5p^5q^5
 =\frac{10!}{5!5!}\frac1{2^5}\frac1{2^5}\approx0{,}246.
\]
\parttitle{2. Не более пяти орлов}
\[
 P(0\le K\le5)=\sum_{k=0}^5C_{10}^kp^kq^{10-k}
 =\bigl(C_{10}^0+C_{10}^1+\cdots+C_{10}^5\bigr)\left(\frac12\right)^{10}.
\]
\parttitle{3. Хотя бы один орёл}
Противоположное событие --- ни одного орла:
\[
 P_{10}(0)=C_{10}^0p^0q^{10}=\left(\frac12\right)^{10},\qquad
 P(K\ge1)=1-\left(\frac12\right)^{10}.
\]

\sourcepage{IMG_7649.jpeg}
\section*{Пять игральных костей}
Бросают пять игральных костей. Найти вероятность того, что на трёх
костях выпадет пятёрка.
Это $n=5$ независимых испытаний: успех --- пятёрка,
$p=1/6$, $q=5/6$.
\[
 P_5(3)=C_5^3p^3q^2
 =\frac{5!}{3!2!}\left(\frac16\right)^3\left(\frac56\right)^2
 \approx0{,}032.
\]

\section*{Рождение девочек}
Вероятность рождения мальчика $0{,}515$, девочки --- $0{,}485$.
В семье шестеро детей. Найти вероятность того, что девочек не более двух.
Полагаем $n=6$, $p=0{,}485$, $q=0{,}515$:
\begin{align*}
 P(K\le2)&=P_6(0)+P_6(1)+P_6(2)\\
 &=C_6^0q^6+C_6^1pq^5+C_6^2p^2q^4\approx0{,}372.
\end{align*}

\sourcepage{IMG_7650.jpeg}
\section*{Лотерейные билеты}
Требуется определить, сколько билетов нужно купить, чтобы вероятность
хотя бы одного выигрыша была не меньше $0{,}9$.
В записанном решении используется вероятность проигрыша по одному билету $q=0{,}99$.
\editorial{Число в первой строке условия записано неполно;
значение $q=0{,}99$ взято из дальнейшего вычисления.}
$B_n$ --- среди $n$ билетов хотя бы один выигрышный;
$\overline B_n$ --- все билеты проигрышные.
Для независимых испытаний
\[
 P(\overline B_n)=P_n(0)=q^n,\qquad
 P(B_n)=1-P(\overline B_n).
\]
Требование $P(B_n)\ge0{,}9$ даёт
\[
 0{,}99^n\le0{,}1,\qquad n\ln0{,}99\le\ln0{,}1,\qquad \ln0{,}99<0.
\]
\editorial{Внизу страницы начата, но не дописана дробь для нижней границы $n$.
Числовой ответ в исходнике отсутствует.}

\sourcepage{IMG_7651.jpeg}
\section*{Формула Пуассона}
В конспекте: $n$ велико, $p$ мало, $\lambda=np$; рядом указан ориентир
$\lambda\le10$.
\[
 \boxed{P_n(k)\approx\frac{\lambda^k e^{-\lambda}}{k!}.}
\]
\section*{Бракованные детали в коробке}
Вероятность брака $p=0{,}015$. Детали упаковывают по $100$ штук.
Найти вероятность того, что в коробке нет бракованных деталей.
\[
 n=100,\qquad p=0{,}015,\qquad q=0{,}985.
\]
Событие $B$ --- ни одной бракованной детали:
\[
 P(B)=P_{100}(0)=C_{100}^0p^0q^{100}.
\]
\[
 \lambda=np=100\cdot0{,}015=1{,}5,\qquad
 P(B)\approx\frac{1{,}5^0e^{-1{,}5}}{0!}=e^{-1{,}5}.
\]

\sourcepage{IMG_7652.jpeg}
\section*{Задача 3.41. Булочки с изюмом}
Для выпечки $1000$ булочек требуется $10$ кг изюма.
Наудачу выбирается булочка: рассматривается вероятность того,
что в ней не окажется изюма. В следующем вопросе рассматриваются
пять выбранных булочек: ровно две без изюма, остальные --- с изюмом.
\editorial{Конец второго вопроса на фотографии тетради обрезан.
Его смысл и решение уточнены по фотографии доски IMG\_7614.jpeg.}
В тетради и на доске используются параметры
\[
 n=1000,\qquad \lambda=\frac{10000}{1000}=10.
\]
\editorial{Параметры сохранены из исходников. Переход от массы изюма
к числу изюминок в них не пояснён.}

\sourcepage{IMG_7614.jpeg (фотография доски)}
\parttitle{1. В выбранной булочке нет изюма}
На доске $A$ --- изюминка попала в рассматриваемую булочку,
$\overline A$ --- не попала. Пусть $B$ --- в булочке нет изюма.
По формуле Пуассона
\[
 P(B)=P_{1000}(0)\approx\frac{10^0e^{-10}}{0!}=e^{-10}.
\]

\parttitle{2. Среди пяти булочек ровно две без изюма}
Обозначим это событие через $C$. Для пяти испытаний используется
формула Бернулли: успех означает, что булочка не содержит изюма.
Подставляя полученное приближение $P(B)\approx e^{-10}$, получаем
\[
 \boxed{P(C)=P_5(2)\approx C_5^2\bigl(e^{-10}\bigr)^2
 \bigl(1-e^{-10}\bigr)^3.}
\]
\end{document}
